% SHARD-ID: AQM-27-STATE-PRESHEAF-QUASILOCAL-ALGEBRA
% SHARD-TITLE: State Presheaf of the Quasi-Local Algebra
% SHARD-SUMMARY: Defines local quantum state spaces for the Zariski observable net and proves they form a presheaf.
% SHARD-SUMMARY: Proves that the state presheaf is not generally a sheaf, by explicit finite marginal counterexamples.
% SHARD-SUMMARY: Identifies global states on the algebraic quasi-local algebra with compatible inverse-limit families of local states.
% SHARD-KEYWORDS: quantum state, density matrix, presheaf, sheaf failure, marginal, partial trace, quasi-local algebra

\section{State Presheaf of the Quasi-Local Algebra}
\label{sec:state-presheaf-quasilocal-algebra}

\subsection{Status and Source Anchors}

This shard adds the state side of AQM-26.  Derezinski records that normal
states on \(B(\mathcal H)\) are given by trace-class functionals
\(\psi(A)=\operatorname{Tr}(\gamma A)\), with positive trace-one operators
\(\gamma\) called density matrices, in
\path{references/canonical_fields/derezinski_math_ph_0511030_source/derezinski.tex},
lines 3141--3149.  Gottesman records the finite quantum-information
convention that density matrices are positive trace-one operators and that a
subsystem density matrix is obtained by tracing out the remaining degrees of
freedom in
\path{references/symplectic_qecc/gottesman_9705052_source/Thesis.tex},
lines 404--411.

\subsection{Lamport Dependency Skeleton}

\[
\begin{array}{rcl}
D1&:&\operatorname{State}(A)=
      \{\omega:A\to\mathbb C\mid \omega(1)=1,\
      \omega(a^*a)\geq0\}.\\
D2&:&\text{For }A=\operatorname{End}(\mathcal H),\
      \omega(a)=\operatorname{Tr}(\rho a),\quad
      \rho\geq0,\ \operatorname{Tr}\rho=1.\\
D3&:&\mathsf S(U)=\operatorname{State}(\mathcal A(U)).\\
L1&:&U\subset V\Rightarrow
      \mathsf S(V)\to\mathsf S(U),\ \omega\mapsto
      \omega\circ\iota_{U,V},\text{ is a state restriction.}\\
T1&:&\mathsf S\text{ is a presheaf of convex sets.}\\
T2&:&\mathsf S\text{ is not generally separated, hence not a sheaf.}\\
T3&:&\text{Compatible overlapping marginals need not glue to a global state.}\\
T4&:&\operatorname{State}(\mathcal A_{\rm qloc}(X))
      \cong\varprojlim_U\operatorname{State}(\mathcal A(U))
      \text{ for the algebraic filtered colimit.}
\end{array}
\]

\subsection{State Spaces D1--D3}

Let \(A\) be a unital \(*\)-algebra.  A state on \(A\) is a linear functional
\(\omega:A\to\mathbb C\) such that
\[
  \omega(1)=1,\qquad \omega(a^*a)\geq0\quad\text{for all }a\in A.
\]
For a finite-dimensional Hilbert space \(\mathcal H\) and
\(A=\operatorname{End}(\mathcal H)\), we use the density-matrix
representation
\[
  \omega_\rho(a)=\operatorname{Tr}(\rho a),
  \qquad \rho\geq0,\quad \operatorname{Tr}\rho=1.
\]
The state space is convex: if \(\omega_0,\omega_1\) are states and
\(0\leq t\leq1\), then \(t\omega_0+(1-t)\omega_1\) is again normalized and
positive.

For the open-local observable net \(U\mapsto\mathcal A(U)\), define
\[
  \mathsf S(U)=\operatorname{State}(\mathcal A(U)).
\]

\subsection{Restriction Is a Presheaf L1--T1}

\paragraph{Lemma \(L1\).}
If \(U\subset V\), then
\[
  \rho^V_U:\mathsf S(V)\longrightarrow\mathsf S(U),
  \qquad
  \rho^V_U(\omega)=\omega\circ\iota_{U,V}
\]
is a well-defined affine map.

\emph{Proof.}
The inclusion \(\iota_{U,V}\) is a unital \(*\)-homomorphism by AQM-24.  Thus
\[
  \rho^V_U(\omega)(1)=\omega(\iota_{U,V}(1))=\omega(1)=1.
\]
For \(a\in\mathcal A(U)\),
\[
  \rho^V_U(\omega)(a^*a)
  =
  \omega(\iota_{U,V}(a^*a))
  =
  \omega(\iota_{U,V}(a)^*\iota_{U,V}(a))\geq0.
\]
Hence \(\rho^V_U(\omega)\) is a state.  Affineness follows from linearity of
composition with \(\iota_{U,V}\).

When \(V=U\sqcup W\), the same restriction is the density-matrix partial
trace.  If \(\omega_\rho\) is a state on
\(\mathcal A(U)\otimes\mathcal A(W)\), then the restricted state on
\(\mathcal A(U)\) is represented by
\[
  \rho_U=\operatorname{Tr}_W(\rho),
\]
because \(\operatorname{Tr}(\rho(a\otimes1_W))
=\operatorname{Tr}(\rho_Ua)\).

\paragraph{Theorem \(T1\).}
\(\mathsf S\) is a presheaf of convex sets on \(X\).

\emph{Proof.}
For \(U\subset V\subset W\),
\[
  \rho^W_U(\omega)
  =
  \omega\circ\iota_{U,W}
  =
  \omega\circ\iota_{V,W}\circ\iota_{U,V}
  =
  \rho^V_U(\rho^W_V(\omega)),
\]
using functoriality of \(\iota\) from AQM-24.  Also
\(\rho^U_U=\operatorname{id}\).  The restriction maps are affine by Lemma
\(L1\), so this is a contravariant functor from opens to convex sets.

\subsection{Failure of Separatedness T2}

\paragraph{Theorem \(T2\).}
\(\mathsf S\) is not generally a sheaf.  In fact it is not generally
separated: two global states can have the same restrictions to an open cover.

\emph{Proof.}
Use \(X=\operatorname{Spec}(\mathbb Z/6\mathbb Z)\), with the two singleton
opens
\[
  U=\{(2)\},\qquad V=\{(3)\},\qquad X=U\sqcup V.
\]
Then \(\mathcal A(X)=M_2\otimes M_3\).  Choose bases
\(|0\rangle,|1\rangle\) of \(\mathbb C^2\) and
\(|0\rangle,|1\rangle,|2\rangle\) of \(\mathbb C^3\), and set
\[
  |\psi\rangle=\frac{|0,0\rangle+|1,1\rangle}{\sqrt2},
  \qquad
  \rho=|\psi\rangle\langle\psi|.
\]
Its restrictions are
\[
  \rho_U=\frac{I_2}{2},\qquad
  \rho_V=\frac{|0\rangle\langle0|+|1\rangle\langle1|}{2}.
\]
Now define the product density matrix
\[
  \sigma=\rho_U\otimes\rho_V.
\]
The density matrices \(\rho\) and \(\sigma\) have the same restrictions to
\(U\) and to \(V\) by construction.  They are not equal: \(\rho\) has the
off-diagonal matrix entry \(|0,0\rangle\langle1,1|/2\), while \(\sigma\) has
no such entry.  Thus the two global states agree on the cover
\(\{U,V\}\) but differ globally.  The uniqueness part of the sheaf condition
fails.

\subsection{Failure of Existence for Overlapping Marginals T3}

\paragraph{Theorem \(T3\).}
Even when local states agree on an overlap, a global extending state need not
exist.

\emph{Proof.}
Consider a discrete three-point finite spectrum with qubit factors labelled
\(A,B,C\).  For example, one may take
\(\operatorname{Spec}(\mathbb F_2\times\mathbb F_2\times\mathbb F_2)\).
Let
\[
  U=\{A,B\},\qquad V=\{B,C\},\qquad U\cap V=\{B\}.
\]
On \(AB\), take the Bell state
\[
  \rho_{AB}=|\Phi_{AB}\rangle\langle\Phi_{AB}|,\qquad
  |\Phi_{AB}\rangle=\frac{|0,0\rangle+|1,1\rangle}{\sqrt2}.
\]
On \(BC\), take
\[
  \rho_{BC}=|\Phi_{BC}\rangle\langle\Phi_{BC}|,\qquad
  |\Phi_{BC}\rangle=\frac{|0,0\rangle+|1,1\rangle}{\sqrt2}.
\]
Both restrict to \(I_2/2\) on the overlap \(B\).

Assume a density matrix \(\theta_{ABC}\) extends both.  Since
\(\operatorname{Tr}_C\theta_{ABC}=\rho_{AB}\) is the rank-one projection
\(P_{AB}=|\Phi_{AB}\rangle\langle\Phi_{AB}|\), positivity forces
\(\theta_{ABC}\) to be supported on
\(\operatorname{ran}(P_{AB})\otimes\mathbb C^2_C\).  Indeed,
\[
  \operatorname{Tr}\bigl(((1-P_{AB})\otimes1_C)\theta_{ABC}\bigr)
  =
  \operatorname{Tr}((1-P_{AB})\rho_{AB})=0,
\]
and a positive operator with zero trace is zero on that positive part.
Therefore
\[
  \theta_{ABC}=P_{AB}\otimes\tau_C
\]
for some density matrix \(\tau_C\).  Its \(BC\)-marginal is then
\[
  \operatorname{Tr}_A(\theta_{ABC})
  =
  \frac{I_2}{2}\otimes\tau_C,
\]
which is a product state.  It cannot equal the pure entangled Bell state
\(\rho_{BC}\).  Hence the compatible pair \((\rho_{AB},\rho_{BC})\) does not
glue to a state on \(ABC\).

\subsection{Global States as an Inverse Limit T4}

Let \(\mathsf K(X)\) be the directed poset of quasi-compact opens used in
AQM-25, and let
\[
  \mathcal A_{\rm qloc}(X)=\varinjlim_{U\in\mathsf K(X)}\mathcal A(U)
\]
be the algebraic quasi-local \(*\)-algebra.  Since the structure maps are
injective, we may regard \(\mathcal A_{\rm qloc}(X)\) as the algebraic union
of the local algebras.

\paragraph{Theorem \(T4\).}
There is a natural affine bijection
\[
  \operatorname{State}(\mathcal A_{\rm qloc}(X))
  \cong
  \varprojlim_{U\in\mathsf K(X)}
  \operatorname{State}(\mathcal A(U)).
\]

\emph{Proof.}
Given a state \(\omega\) on \(\mathcal A_{\rm qloc}(X)\), restrict it to each
\(\mathcal A(U)\).  These restrictions are compatible by functoriality of
restriction, so they define an element of the inverse limit.

Conversely, let \((\omega_U)_U\) be a compatible family of local states.  If
\(a\in\mathcal A_{\rm qloc}(X)\), choose \(U\) and \(a_U\in\mathcal A(U)\)
representing it, and define
\[
  \omega(a)=\omega_U(a_U).
\]
This is independent of the representative: if \(a_U\) and \(a_V\) represent
the same quasi-local element, choose \(W\supset U,V\); their images in
\(\mathcal A(W)\) agree, and compatibility gives the same value.  Linearity
follows after moving two representatives to a common \(W\).  Normalization is
\(\omega(1)=1\).  Positivity holds because if
\(a\in\mathcal A(U)\), then \(a^*a\in\mathcal A(U)\) and
\[
  \omega(a^*a)=\omega_U(a^*a)\geq0.
\]
The two constructions are inverse to each other.

\subsection{Conclusion}

The state side reverses the observable-net arrows:
\[
  U\subset V\quad\Rightarrow\quad
  \operatorname{State}(\mathcal A(V))
  \longrightarrow
  \operatorname{State}(\mathcal A(U)).
\]
Thus local quantum states form a presheaf of convex sets.  They do not form a
sheaf in general: uniqueness fails because marginals do not determine
correlations, and existence fails because compatible overlapping marginals may
not have a joint quantum state.  The correct quasi-local replacement is the
inverse-limit statement of Theorem \(T4\), not an ordinary sheaf gluing axiom.
